% page 419 of the facsimile (image p-418 of the scan)
% block 162.6 x 237.7 mm ; left margin 39.4 mm
\pgc{17.780mm}{0.000mm}{
\l{\cel{52}411.}
\l{}
\l{\sou{ornitologio}. Historio naturala pri la ucelo. -\cel{2}DEFIRS.}
\l{}
\l{\sou{ornitoptero}. Papiliono belega, la maxim granda ek le jornala,}
\l{\cel{3}qua vivas en la foresti humida, e flugas tre alte. - DEFIRS.}
\l{}
\l{\sou{orografio}. Deskripto di la monti.DEFIRS.}
\l{}
\l{\sou{orpimento}. Sulfido flava di arseno, uzata kom.farbo e por sen-}
\l{\cel{3}piligar la feli.\cel{2}Formulo kemiala : As2\cel{2}S3 .\cel{2}DEFIS.}
\l{}
\l{\sou{orselo}. (\sou{bot.)}\cel{2}Likeno qua kreskas precipue en Madagaskar-}
\l{\cel{3}insulo, e donas bela farbo violea. - L.\cel{1}\sou{roccella tinctoria}.}
\l{}
\l{\sou{orta}. (\sou{geom.}) (\sou{pri angulo}). Di qua la aperturo kontenas 90}
\l{\cel{3}gradi di la cirklo-kurvo. -}
\l{}
\l{\sou{ortocentro}. (\sou{geom.}) Nomo di la konvergo-punto di la tri altesi}
\l{\cel{3}di triangulo. -}
\l{}
\l{ortodoxa. (\sou{teol. kristana)}. Konforma a la doktrino korekta. -}
\l{\cel{3}\sou{(jokoze}). Qua esas to quon lu devas esar.-DEFIRS.}
\l{}
\l{ortogenezo. (\sou{biol.)}\cel{2}Vario spontana, pro kauzo interna, qua}
\l{\cel{3}pulsas la organismi a punto determinita, sen la interveno}
\l{\cel{3}da adapto. -}
\l{}
\l{ortogono. (\sou{geom.}) Qua formacas angulo orta - qua esas per-}
\l{\cel{3}pendikla. -}
\l{}
\l{ortografio. (\sou{gram.}) Maniero skribar linguo korekte. Maniero}
\l{\cel{3}skribar la vorti. - - DEFIRS.}
\l{}
\l{ortografika. (\sou{matem.)}\cel{2}Qua reprezentas objekto sur plano,lua}
\l{\cel{3}omna punti esante projektata perpendikle sur la plano.- DEFI.}
\l{}
\l{ortokrona. (\sou{fotogr.)}\cel{2}(\sou{pri la plaki o filmi fotografala}).}
\l{\cel{3}Di qua la sentiveso pri la kolori diversa di la spektro}
\l{\cel{3}igesis, per la ajunto di korpi kemiala, harmoniar kun la}
\l{\cel{3}sentiveso di la hom-okulo. -}
\l{}
\l{ortolano. (\sou{zool.}) Mikra ucelo migranta, qua vivas en Mezeuropa,}
\l{\cel{3}ed en Sudeuropa, di qua la karno (manjebla) esas delikata.}
\l{\cel{3}- L. emberiza hortulana. - DEFIS.}
\l{}
\l{ortopedio. La arto preventar o pokope supresar la deformuri}
\l{\cel{3}di la korpo (precipue che la pueri), per aparato o kuracado.}
\l{\cel{3}- DEFIRS.}
\l{}
\l{\cel{1}\sou{ortoptero}. (\sou{zool.}) Insekto qua havas quar ali, di qui la du}
\l{\cel{3}infra esas faldita longeso-sinse. - DEFIS.}
\l{}
\l{\cel{1}\sou{ortoptika}.\cel{2}(\sou{matem.}) (\sou{pri kurvo}) Loko geometriala di la punti}
\l{\cel{4}de ube onu vidas kurvo plana donata ortangule. -}
}
